% Generated draft: requires mathematical and source review.
\documentclass[11pt]{article}
\usepackage[margin=25mm]{geometry}
\usepackage{fontspec}
\setmainfont{Latin Modern Roman}
\usepackage{amsmath,amssymb,mathtools,unicode-math}
\setmathfont{Latin Modern Math}
\usepackage{enumitem,needspace,xurl,hyperref}
\setlength{\parindent}{0pt}
\setlength{\parskip}{5pt}
\setlength{\emergencystretch}{3em}
\newcommand{\N}{\mathbb{N}}
\newcommand{\Z}{\mathbb{Z}}
\newcommand{\Q}{\mathbb{Q}}
\newcommand{\R}{\mathbb{R}}
\newcommand{\C}{\mathbb{C}}
\newcommand{\F}{\mathbb{F}}
\newcommand{\A}{\mathbb{A}}
\newcommand{\PP}{\mathbb{P}}
\newcommand{\E}{\mathbb{E}}
\newcommand{\eps}{\varepsilon}
\newcommand{\dd}{\,\mathrm{d}}
\begin{document}
\section*{UnsolvedMath ID: 30004352}
\textbf{Source problem code:} catalogue entry
\subsection{Definitions and mathematical statement}
Let \(n\in\Z_{\geq 1}\), and let \(G=(V,E)\) be a graph. A matching is a set of edges whose sets of incident vertices are pairwise disjoint; its size is its number of edges. Let \((M_i)_{i=1}^{n}\) be an indexed family of matchings in \(G\), with \(|M_i|=n\) for every \(i\in\{1,\ldots,n\}\). The matchings are not required to be disjoint from one another. A rainbow matching for this family is a matching \(R\subseteq E\) together with an injective map \(c:R\to\{1,\ldots,n\}\) such that \(e\in M_{c(e)}\) for every \(e\in R\). Thus an edge belonging to several matchings is assigned just one index in a rainbow matching.

\textbf{Question.} Does every such family in every graph contain a rainbow matching of size \(n-1\)? Equivalently, must there exist \(I\subseteq\{1,\ldots,n\}\) with \(|I|=n-1\) and edges \(e_i\in M_i\) for \(i\in I\) whose sets of incident vertices are pairwise disjoint? For \(n=1\), the required rainbow matching is empty.

\textbf{Unresolved qualifications.} The source imposes no bipartiteness hypothesis or explicit finiteness restriction on \(G\). It does not explicitly specify graph conventions concerning directed edges, loops, or parallel edges; the quantification over graphs retains this source limitation and requires review. The family is indexed to make the rainbow condition precise; the source does not explicitly discuss repeated identical matchings.
\subsection{Short English statement}
Given \(n\) matchings containing \(n\) edges each in a graph, can one choose \(n-1\) pairwise vertex-disjoint edges and assign them to distinct matchings that contain them? The supplied curation identifies this as the Aharoni--Berger rainbow-matching conjecture and labels it open; that attribution and status have not been independently verified.
\subsection{Sources}
\begin{itemize}
\item Local UnsolvedMath registry: supplied complete record for permanent ID \(30004352\), titled \emph{Rainbow Matchings Across Families of Matchings}, in the Oberwolfach Reports Open Problems collection. Exact source problem code: \texttt{catalogue entry}. The record gives proposed year \(2020\), supplies no individual proposer, and labels the statement \texttt{corrected\_verified}; these are supplied metadata, not independent mathematical verification.
\item Supplied report citation: \emph{Combinatorics (2021)}, Oberwolfach Reports, DOI 10.4171/owr/2020/1, pp.\ 6--89. \url{https://doi.org/10.4171/owr/2020/1}. The background describes these reports as documenting mathematical workshops and their open-problem discussions. The supplied citation has not been consulted independently.
\item Aharoni--Berger--Chudnovsky--Zerbib, ``rainbow matchings'' (supplied reference label). \url{https://doi.org/10.37236/10173}.
\item Bowtell--Freschi--Kronenberg--Yan, ``improved bounds'' (supplied reference label). \url{https://arxiv.org/abs/2501.03216}.
\item Supplied literature assessment, checked \texttt{2026-08-22}: the Aharoni--Berger conjecture for \(n\) matchings of size \(n\) in an arbitrary graph is labeled open; the assessment mentions lower bounds such as \(2n/3-O(1)\) and stronger bounds in restricted graph classes. This is supplied dated curation derived from the OWR material, not an independently verified literature assessment. The source itself requests independent verification of time-sensitive status and cited claims; it supplies no precise hypotheses, constants, or theorem references for these bounds.
\end{itemize}
\end{document}
